\section{真空システムとチャンバー設計}

\subsection{差動排気システム (Differential Pumping)}
本装置の心臓部であるMCP（F2222-21P）の動作推奨真空度は $1.3 \times 10^{-4}\,\mathrm{Pa}$ ($1.0 \times 10^{-6}\,\mathrm{Torr}$) 以下である。
しかし、TS-6本体の実験中の真空度は $8.0 \times 10^{-4}\,\mathrm{Pa}$ ($~10^{-5}\,\mathrm{Torr}$) 程度まで悪化することがあり、そのまま接続するとMCPが放電・破壊するリスクがある\cite{Takeda2023}。

そのため、TS-6本体と計測器の間で圧力差を作る「差動排気」システムを採用している。
この成否は、以下の2点のバランスで決まる。
\begin{itemize}
    \item コンダクタンスの制限（ピンホール）:
    プラズマとの接続部を $\phi 1\,\mathrm{mm}$ の微小ピンホールのみに制限することで、ガスの流入量を極限まで抑える\cite{Xiang2021}。
    \item 独立排気（専用TMP）:
    流入した僅かなガスを、MCPチャンバーに直結した小型ターボ分子ポンプ（TMP）で直ちに排気する\cite{Takeda2023}。
\end{itemize}

\subsection{チャンバー設計の詳細と進化}
\subsubsection{初期型（独立型）から4視点一体型へ}
初期の設計では1視点につき1つの真空容器を使用していたが、視点数を増やす際に配管が複雑化し、真空リークのリスクも増大した。
これを解決するため、修論第5章にある通り、「4つのMCPを単一の筐体に収める（4-in-1）」設計へと変更した\cite{Takeda2023}。

\begin{itemize}
    \item 材質: SUS304（ステンレス鋼）。外部磁場（10-100 kHz）の侵入を防ぎ、ノイズを遮蔽するために 板厚 5mm を確保している\cite{Xiang2021}。
    \item 構造: 直方体の筐体内部に4つのMCPマウントを配置し、単一のNWフランジ（ポート）から排気と電圧供給ケーブルの導入を一括で行う構造とした\cite{Takeda2023}。
    \item 寸法: MCP有効径（$\phi 17\,\mathrm{mm}$）とマウントサイズを考慮し、内部干渉しない最小限のサイズで設計されている。
\end{itemize}

\subsection{ピンホールの「二重の意図」}
本装置におけるピンホール（$\phi 1\,\mathrm{mm}$）は、単なる光学的な穴ではない。
\begin{enumerate}
    \item 光学素子として: ピンホールカメラの原理により像を結ぶための開口部。穴径が小さいほど分解能は上がるが、光量は減る（トレードオフ）。
    \item 真空隔壁として: オリフィス（Orifice）として機能し、TS-6本体（低真空）とMCP室（高真空）の圧力差を維持する「抵抗」の役割を果たす。この穴が大きすぎると差動排気が破綻し、MCPが破壊される。
\end{enumerate}

\subsection{ベーキング（加熱）時の「絶対」注意事項}
到達真空度を向上させるため、リボンヒーターによる真空容器の加熱（ベーキング）を行う場合がある\cite{Takeda2023}が、以下の点に注意すること。

\begin{itemize}
    \item Oリングの耐熱温度:
    真空シールに使用しているOリング（Viton等）の耐熱温度（通常 $150^\circ\mathrm{C}$ 程度）を超えないように温度管理を行うこと。過熱すると硬化・変形し、真空リークの原因となる。
    \item 3Dプリンタ製パーツの溶解:
    暗室ボックスやカメラ固定具などにPLA/ABS等の3Dプリンタ製パーツを使用している場合、ベーキングの熱が伝わると容易に軟化・変形する。加熱時は必ず周囲の樹脂パーツを取り外すか、熱絶縁を徹底すること。
    \item MCPへの熱影響:
    MCP自体も急激な温度変化に弱いため、昇温・降温は時間をかけてゆっくり行う（$10^\circ\mathrm{C}/\mathrm{min}$ 以下推奨）。
\end{itemize}